Quadratic unconstrained binary optimization (QUBO) lets annealing hardware attack scheduling, routing, location
and many other combinatorial problems, but turning a real problem into a solved QUBO still takes an expert: the
problem must be pinned down, real data must be found and cleaned, the model must be written and checked, and
penalty weights and solver settings must be tuned on the hardware. Large language models
(LLMs) can draft QUBO formulations from natural language, but formulation correctness alone does not establish
that a workflow has acquired the intended data or delivered a usable decision.
We present a human-in-the-loop agent system that goes from a one-sentence problem description to a solution on
an annealer, built entirely on a locally served open model. The system clarifies the problem (with sourced web
lookup), lets the model choose and write the formulation under problem-independent structural checks, finds and
preprocesses public data through a typed sets-and-parameters interface, records provenance, and submits the QUBO to
interchangeable backends: simulated annealing, a GPU annealer, Fujitsu's Digital Annealer and D-Wave (direct QPU
and the Leap Hybrid service). The experiments reported here use simulated annealing, the GPU annealer and the
D-Wave Leap Hybrid service; the Digital Annealer rejected our submission at its monthly account limit. An outer
recovery loop checks data and model failures before any tuning, preserves model--data versions, and
distinguishes usable, limited and unavailable solutions. Penalty tuning is a secondary option for improving feasibility or quality against a baseline; the UI
starts with the default policy.
Separate checker roles compare prepared data with raw-table aggregates and model decisions with the confirmed
requirements. These checks provide repair evidence rather than correctness guarantees. Optional retrieval
supplies general modelling patterns. A person confirms the specification, approves data assumptions and reviews
the result; unresolved failures are returned with a diagnosis. A new real-data ablation protocol separates
pipeline structure, checking, repair and pattern retrieval; repeated evaluation under a frozen version remains
open.
We also show that the standard way of verifying LLM-written QUBOs is unreliable: in an audit of 33 failures of
an existing pipeline, at least 7 of 12 ``shape'' failures were caused by wrong LLM-written test cases rather than
wrong models, and one benchmark's reference answer contradicts itself. We propose problem-level verification,
which decodes every ground state and checks it with an independent checker anchored on the human-stated
optimum. Under it, the pipeline solves 70.3\% (69--72\%) of 100 benchmarks over three runs, while 27\% of the
benchmarks change outcome between identical runs. On a 12-category suite the local system produced a valid model
for all 36 runs, and the GPU annealer reached the exact optimum in 7 of the 9 categories with a known optimum.
